\scanpage{86}
\mspage{7}
% archive foliation: 85

groupoïdes envisagé est représentable par le schéma
\[
(37)\qquad U_{03}\times\unsure{\boldsymbol{\mu}|\mu}_4^*
\]
Le fibré $\Delta_E$ correspondant à la courbe elliptique épinglée (par Legendre-Torelli) universelle sur la base $U_{03}\times\unsure{\boldsymbol{\mu}|\mu}_4^*$, est isomorphe canoniquement à
\[
\deleted{\underline{O}_{U_{0,3}}(1)}\qquad\deleted{\unsure{T}_0=\unsure{\underline{O}^{\sharp}}\,\underline{O}_{U_{0,3}}(1)^{\unsure{\sharp}}\times\unsure{\boldsymbol{\mu}|\mu}}
\]
\[
(38)\qquad\mathbb{T}_0=\underline{O}_{U_{0,3}\times\unsure{\boldsymbol{\mu}|\mu}_4^*}(1)^{\unsure{\sharp}}
\]
où ici, puisqu'on est \deleted{sur}\inserted{au dessus} $\unsure{\boldsymbol{\mu}|\mu}_4^*$, le fibré en carrés $Q_{0S}$ $(S=U_{03}\times\unsure{\boldsymbol{\mu}|\mu}_4^*)$ est trivial, et on a un iso. can.
\[
(39)\qquad\mathbb{T}_0\simeq\underline{O}_{U_{0,3}}(1)
\]
\deleted{Le groupe modulaire pour \unsure{l'étude}}\inserted{Mais il faut expliciter le ``groupe} modulaire'' naturel pour cette situation des \deleted{et le groupe} structures épinglées, et examiner \uline{son} action sur $U_{03}\times\unsure{\boldsymbol{\mu}|\mu}_4^*$ ainsi que sur $\mathbb{T}_0$.
